\section{Discussion}

The current evidence supports a focused claim: multi-relation role graph learning improves kill-chain recovery under strict intermittent sensing and relay-node communication failure in a constrained 3DOF heterogeneous UAV interception task. The fixed-budget five-seed evaluation strengthens this claim because all methods use the same checkpoint rule, matched test episodes, and seed-aware statistical analysis. The full multi-relation method improves recovery over both no-graph and single-graph baselines while maintaining zero test collisions.

The strongest mechanism evidence is the role-pair-conditioned message-gating ablation. Removing the role-pair gate reduces recovery and increases timeout with a seed-aware confidence interval that separates in favor of the full method. This supports the design choice that messages should depend not only on graph connectivity but also on the sender role, receiver role, and relation context. The task-support relation ablation is directionally useful but less statistically clean: it lowers mean recovery, but the seed-aware interval crosses zero. It should therefore be interpreted as supportive evidence rather than as an isolated decisive proof.

The failure-aligned mechanism curves show that the full method does not simply avoid failure by preserving full communication connectivity. Connectivity drops sharply after relay failure for all methods. The advantage instead appears in higher target tracking and earlier recovery accumulation after the failure window begins. This is consistent with the intended role of multi-relation message passing: preserve and reconstruct task-relevant information under partial sensing and degraded communication rather than assuming a continuously connected team graph.

There are also important limitations. First, the current paper-facing package uses a fixed checkpoint at update 60. This avoids noisy validation selection and keeps all methods under a matched training budget, but it should not be mixed with validation-selected results without explicitly stating the checkpoint-selection protocol. Second, the three-seed no-curriculum diagnostic did not show an independent topology-curriculum advantage. Therefore, the current claim should emphasize graph and message mechanisms, while topology curriculum should be described as a training protocol rather than as a primary contribution. Third, the environment remains a 3DOF mechanism study. It does not yet constitute a complete 4v2 red-blue air-combat system with online missile launch, high-fidelity radar, human-UAV teaming, or 6DOF JSBSim validation.

These limitations define the next development path. The current result package is suitable as a strong 3v1 mechanism foundation. To target a higher-tier submission, future work should add carefully controlled scenario-depth experiments, such as mild maneuvering targets, jammer or escort interference, and limited 6DOF replay validation, without diluting the current evidence chain.
